\chapter{Layout Localization Pipeline}\label{record-049}

\noindent{\footnotesize\color{BookMuted}\path{docs/strategy/layout-localization-pipeline.md}\par}\smallskip

\emph{Strategy note for how Bijou turns \texttt{L\allowbreak{}a\allowbreak{}y\allowbreak{}o\allowbreak{}u\allowbreak{}t\allowbreak{}N\allowbreak{}o\allowbreak{}d\allowbreak{}e} trees into local render roots}

\section{Why this exists}\label{record-049-why-this-exists}

The \texttt{L\allowbreak{}a\allowbreak{}y\allowbreak{}o\allowbreak{}u\allowbreak{}t\allowbreak{}N\allowbreak{}o\allowbreak{}d\allowbreak{}e} path in \texttt{@\allowbreak{}f\allowbreak{}l\allowbreak{}y\allowbreak{}i\allowbreak{}n\allowbreak{}g\allowbreak{}r\allowbreak{}o\allowbreak{}b\allowbreak{}o\allowbreak{}t\allowbreak{}s\allowbreak{}/\allowbreak{}b\allowbreak{}i\allowbreak{}j\allowbreak{}o\allowbreak{}u\allowbreak{}-\allowbreak{}t\allowbreak{}u\allowbreak{}i} is easy to use from the app side, but the runtime ownership model is easy to misunderstand.

The important distinction is:

\begin{itemize}
\tightlist
\item
  \textbf{layout localization} decides what coordinate space the tree owns
\item
  \textbf{paint} only blits already-positioned surfaces into a target surface
\end{itemize}

If those two stages get blurred together, it becomes difficult to reason about:

\begin{itemize}
\tightlist
\item
  negative coordinates and off-origin roots
\item
  where width and height are measured
\item
  when layout is recomputed
\item
  where retained-layout reuse should fit later
\end{itemize}

This note exists to make that path explicit.

\section{The current runtime path}\label{record-049-the-current-runtime-path}

For a TUI render driven by \texttt{r\allowbreak{}u\allowbreak{}n\allowbreak{}(\allowbreak{}.\allowbreak{}.\allowbreak{}.\allowbreak{})}, the flow is:

\begin{enumerate}
\def\labelenumi{\arabic{enumi}.}
\tightlist
\item
  \texttt{a\allowbreak{}p\allowbreak{}p\allowbreak{}.\allowbreak{}v\allowbreak{}i\allowbreak{}e\allowbreak{}w\allowbreak{}(\allowbreak{}m\allowbreak{}o\allowbreak{}d\allowbreak{}e\allowbreak{}l\allowbreak{})} returns either a \texttt{S\allowbreak{}u\allowbreak{}r\allowbreak{}f\allowbreak{}a\allowbreak{}c\allowbreak{}e} or a \texttt{L\allowbreak{}a\allowbreak{}y\allowbreak{}o\allowbreak{}u\allowbreak{}t\allowbreak{}N\allowbreak{}o\allowbreak{}d\allowbreak{}e}.
\item
  \texttt{w\allowbreak{}r\allowbreak{}a\allowbreak{}p\allowbreak{}V\allowbreak{}i\allowbreak{}e\allowbreak{}w\allowbreak{}O\allowbreak{}u\allowbreak{}t\allowbreak{}p\allowbreak{}u\allowbreak{}t\allowbreak{}A\allowbreak{}s\allowbreak{}L\allowbreak{}a\allowbreak{}y\allowbreak{}o\allowbreak{}u\allowbreak{}t\allowbreak{}R\allowbreak{}o\allowbreak{}o\allowbreak{}t\allowbreak{}(\allowbreak{}.\allowbreak{}.\allowbreak{}.\allowbreak{})} converts the view output into a runtime \texttt{l\allowbreak{}a\allowbreak{}y\allowbreak{}o\allowbreak{}u\allowbreak{}t\allowbreak{}R\allowbreak{}o\allowbreak{}o\allowbreak{}t}.
\item
  If the output was a \texttt{L\allowbreak{}a\allowbreak{}y\allowbreak{}o\allowbreak{}u\allowbreak{}t\allowbreak{}N\allowbreak{}o\allowbreak{}d\allowbreak{}e}, \texttt{l\allowbreak{}o\allowbreak{}c\allowbreak{}a\allowbreak{}l\allowbreak{}i\allowbreak{}z\allowbreak{}e\allowbreak{}L\allowbreak{}a\allowbreak{}y\allowbreak{}o\allowbreak{}u\allowbreak{}t\allowbreak{}N\allowbreak{}o\allowbreak{}d\allowbreak{}e\allowbreak{}(\allowbreak{}.\allowbreak{}.\allowbreak{}.\allowbreak{})} shifts the whole tree into a local non-negative coordinate space.
\item
  Layout-stage middleware such as motion and BCSS operate on that localized root.
\item
  The paint middleware blits \texttt{n\allowbreak{}o\allowbreak{}d\allowbreak{}e\allowbreak{}.\allowbreak{}s\allowbreak{}u\allowbreak{}r\allowbreak{}f\allowbreak{}a\allowbreak{}c\allowbreak{}e} values into the target surface at the node's already-localized \texttt{r\allowbreak{}e\allowbreak{}c\allowbreak{}t\allowbreak{}.\allowbreak{}x} / \texttt{r\allowbreak{}e\allowbreak{}c\allowbreak{}t\allowbreak{}.\allowbreak{}y}.
\item
  Diff and output stages serialize the target surface to terminal bytes.
\end{enumerate}

The relevant seams today are:

\begin{itemize}
\tightlist
\item
  \texttt{p\allowbreak{}a\allowbreak{}c\allowbreak{}k\allowbreak{}a\allowbreak{}g\allowbreak{}e\allowbreak{}s\allowbreak{}/\allowbreak{}b\allowbreak{}i\allowbreak{}j\allowbreak{}o\allowbreak{}u\allowbreak{}-\allowbreak{}t\allowbreak{}u\allowbreak{}i\allowbreak{}/\allowbreak{}s\allowbreak{}r\allowbreak{}c\allowbreak{}/\allowbreak{}v\allowbreak{}i\allowbreak{}e\allowbreak{}w\allowbreak{}-\allowbreak{}o\allowbreak{}u\allowbreak{}t\allowbreak{}p\allowbreak{}u\allowbreak{}t\allowbreak{}.\allowbreak{}t\allowbreak{}s}
\item
  \texttt{p\allowbreak{}a\allowbreak{}c\allowbreak{}k\allowbreak{}a\allowbreak{}g\allowbreak{}e\allowbreak{}s\allowbreak{}/\allowbreak{}b\allowbreak{}i\allowbreak{}j\allowbreak{}o\allowbreak{}u\allowbreak{}-\allowbreak{}t\allowbreak{}u\allowbreak{}i\allowbreak{}/\allowbreak{}s\allowbreak{}r\allowbreak{}c\allowbreak{}/\allowbreak{}l\allowbreak{}a\allowbreak{}y\allowbreak{}o\allowbreak{}u\allowbreak{}t\allowbreak{}-\allowbreak{}n\allowbreak{}o\allowbreak{}d\allowbreak{}e\allowbreak{}-\allowbreak{}s\allowbreak{}u\allowbreak{}r\allowbreak{}f\allowbreak{}a\allowbreak{}c\allowbreak{}e\allowbreak{}.\allowbreak{}t\allowbreak{}s}
\item
  \texttt{p\allowbreak{}a\allowbreak{}c\allowbreak{}k\allowbreak{}a\allowbreak{}g\allowbreak{}e\allowbreak{}s\allowbreak{}/\allowbreak{}b\allowbreak{}i\allowbreak{}j\allowbreak{}o\allowbreak{}u\allowbreak{}-\allowbreak{}t\allowbreak{}u\allowbreak{}i\allowbreak{}/\allowbreak{}s\allowbreak{}r\allowbreak{}c\allowbreak{}/\allowbreak{}r\allowbreak{}u\allowbreak{}n\allowbreak{}t\allowbreak{}i\allowbreak{}m\allowbreak{}e\allowbreak{}.\allowbreak{}t\allowbreak{}s}
\item
  \texttt{p\allowbreak{}a\allowbreak{}c\allowbreak{}k\allowbreak{}a\allowbreak{}g\allowbreak{}e\allowbreak{}s\allowbreak{}/\allowbreak{}b\allowbreak{}i\allowbreak{}j\allowbreak{}o\allowbreak{}u\allowbreak{}-\allowbreak{}t\allowbreak{}u\allowbreak{}i\allowbreak{}/\allowbreak{}s\allowbreak{}r\allowbreak{}c\allowbreak{}/\allowbreak{}p\allowbreak{}i\allowbreak{}p\allowbreak{}e\allowbreak{}l\allowbreak{}i\allowbreak{}n\allowbreak{}e\allowbreak{}/\allowbreak{}m\allowbreak{}i\allowbreak{}d\allowbreak{}d\allowbreak{}l\allowbreak{}e\allowbreak{}w\allowbreak{}a\allowbreak{}r\allowbreak{}e\allowbreak{}/\allowbreak{}p\allowbreak{}a\allowbreak{}i\allowbreak{}n\allowbreak{}t\allowbreak{}.\allowbreak{}t\allowbreak{}s}
\end{itemize}

\section{Coordinate spaces}\label{record-049-coordinate-spaces}

There are three distinct coordinate spaces in play.

\subsection{1. Author space}\label{record-049-author-space}

This is the coordinate space returned by \texttt{a\allowbreak{}p\allowbreak{}p\allowbreak{}.\allowbreak{}v\allowbreak{}i\allowbreak{}e\allowbreak{}w\allowbreak{}(\allowbreak{}m\allowbreak{}o\allowbreak{}d\allowbreak{}e\allowbreak{}l\allowbreak{})}.

In author space:

\begin{itemize}
\tightlist
\item
  a root node may start at negative \texttt{x} or \texttt{y}
\item
  siblings may be positioned relative to a parent origin that is not \texttt{(\allowbreak{}0\allowbreak{},\allowbreak{} \allowbreak{}0\allowbreak{})}
\item
  the bounds of the full tree may extend outside the root node's own rect
\end{itemize}

That is legal input. The runtime does not require app authors to hand-normalize it.

\subsection{2. Local root space}\label{record-049-local-root-space}

This is the coordinate space owned by the runtime pipeline.

\texttt{l\allowbreak{}o\allowbreak{}c\allowbreak{}a\allowbreak{}l\allowbreak{}i\allowbreak{}z\allowbreak{}e\allowbreak{}L\allowbreak{}a\allowbreak{}y\allowbreak{}o\allowbreak{}u\allowbreak{}t\allowbreak{}(\allowbreak{}.\allowbreak{}.\allowbreak{}.\allowbreak{})} recursively measures the tree bounds and computes:

\begin{itemize}
\tightlist
\item
  \texttt{w\allowbreak{}i\allowbreak{}d\allowbreak{}t\allowbreak{}h\allowbreak{} \allowbreak{}=\allowbreak{} \allowbreak{}m\allowbreak{}a\allowbreak{}x\allowbreak{}X\allowbreak{} \allowbreak{}-\allowbreak{} \allowbreak{}m\allowbreak{}i\allowbreak{}n\allowbreak{}X}
\item
  \texttt{h\allowbreak{}e\allowbreak{}i\allowbreak{}g\allowbreak{}h\allowbreak{}t\allowbreak{} \allowbreak{}=\allowbreak{} \allowbreak{}m\allowbreak{}a\allowbreak{}x\allowbreak{}Y\allowbreak{} \allowbreak{}-\allowbreak{} \allowbreak{}m\allowbreak{}i\allowbreak{}n\allowbreak{}Y}
\item
  \texttt{d\allowbreak{}x\allowbreak{} \allowbreak{}=\allowbreak{} \allowbreak{}-\allowbreak{}m\allowbreak{}i\allowbreak{}n\allowbreak{}X}
\item
  \texttt{d\allowbreak{}y\allowbreak{} \allowbreak{}=\allowbreak{} \allowbreak{}-\allowbreak{}m\allowbreak{}i\allowbreak{}n\allowbreak{}Y}
\end{itemize}

\texttt{l\allowbreak{}o\allowbreak{}c\allowbreak{}a\allowbreak{}l\allowbreak{}i\allowbreak{}z\allowbreak{}e\allowbreak{}L\allowbreak{}a\allowbreak{}y\allowbreak{}o\allowbreak{}u\allowbreak{}t\allowbreak{}N\allowbreak{}o\allowbreak{}d\allowbreak{}e\allowbreak{}(\allowbreak{}.\allowbreak{}.\allowbreak{}.\allowbreak{})} then translates the whole tree by \texttt{(\allowbreak{}d\allowbreak{}x\allowbreak{},\allowbreak{} \allowbreak{}d\allowbreak{}y\allowbreak{})} so the localized root and every descendant sit in a non-negative local coordinate space.

The important rule is:

\begin{quote}
Paint should never have to rediscover where the root begins.
\end{quote}

By the time the pipeline reaches paint, the root already owns a local origin.

\subsection{3. Target-surface space}\label{record-049-target-surface-space}

This is the actual framebuffer for the current render.

The target surface may be larger than the localized tree because the runtime normalizes to at least the current viewport size:

\begin{itemize}
\tightlist
\item
  \texttt{w\allowbreak{}i\allowbreak{}d\allowbreak{}t\allowbreak{}h\allowbreak{} \allowbreak{}=\allowbreak{} \allowbreak{}m\allowbreak{}a\allowbreak{}x\allowbreak{}(\allowbreak{}v\allowbreak{}i\allowbreak{}e\allowbreak{}w\allowbreak{}p\allowbreak{}o\allowbreak{}r\allowbreak{}t\allowbreak{}.\allowbreak{}w\allowbreak{}i\allowbreak{}d\allowbreak{}t\allowbreak{}h\allowbreak{},\allowbreak{} \allowbreak{}l\allowbreak{}o\allowbreak{}c\allowbreak{}a\allowbreak{}l\allowbreak{}i\allowbreak{}z\allowbreak{}e\allowbreak{}d\allowbreak{}.\allowbreak{}w\allowbreak{}i\allowbreak{}d\allowbreak{}t\allowbreak{}h\allowbreak{})}
\item
  \texttt{h\allowbreak{}e\allowbreak{}i\allowbreak{}g\allowbreak{}h\allowbreak{}t\allowbreak{} \allowbreak{}=\allowbreak{} \allowbreak{}m\allowbreak{}a\allowbreak{}x\allowbreak{}(\allowbreak{}v\allowbreak{}i\allowbreak{}e\allowbreak{}w\allowbreak{}p\allowbreak{}o\allowbreak{}r\allowbreak{}t\allowbreak{}.\allowbreak{}h\allowbreak{}e\allowbreak{}i\allowbreak{}g\allowbreak{}h\allowbreak{}t\allowbreak{},\allowbreak{} \allowbreak{}l\allowbreak{}o\allowbreak{}c\allowbreak{}a\allowbreak{}l\allowbreak{}i\allowbreak{}z\allowbreak{}e\allowbreak{}d\allowbreak{}.\allowbreak{}h\allowbreak{}e\allowbreak{}i\allowbreak{}g\allowbreak{}h\allowbreak{}t\allowbreak{})}
\end{itemize}

That means localization decides the tree's own bounds, while normalization decides the backing framebuffer size.

\section{Localization is not paint}\label{record-049-localization-is-not-paint}

Localization and paint are separate responsibilities.

\subsection{Localization does:}\label{record-049-localization-does}

\begin{itemize}
\tightlist
\item
  recursive bounds measurement
\item
  ownership of root translation
\item
  conversion from author coordinates to local runtime coordinates
\item
  reporting the minimal localized width and height
\end{itemize}

\subsection{Paint does:}\label{record-049-paint-does}

\begin{itemize}
\tightlist
\item
  walk the already-localized tree
\item
  blit \texttt{n\allowbreak{}o\allowbreak{}d\allowbreak{}e\allowbreak{}.\allowbreak{}s\allowbreak{}u\allowbreak{}r\allowbreak{}f\allowbreak{}a\allowbreak{}c\allowbreak{}e} into the target surface
\item
  trust the rects it is given
\end{itemize}

The paint middleware in \texttt{p\allowbreak{}a\allowbreak{}c\allowbreak{}k\allowbreak{}a\allowbreak{}g\allowbreak{}e\allowbreak{}s\allowbreak{}/\allowbreak{}b\allowbreak{}i\allowbreak{}j\allowbreak{}o\allowbreak{}u\allowbreak{}-\allowbreak{}t\allowbreak{}u\allowbreak{}i\allowbreak{}/\allowbreak{}s\allowbreak{}r\allowbreak{}c\allowbreak{}/\allowbreak{}p\allowbreak{}i\allowbreak{}p\allowbreak{}e\allowbreak{}l\allowbreak{}i\allowbreak{}n\allowbreak{}e\allowbreak{}/\allowbreak{}m\allowbreak{}i\allowbreak{}d\allowbreak{}d\allowbreak{}l\allowbreak{}e\allowbreak{}w\allowbreak{}a\allowbreak{}r\allowbreak{}e\allowbreak{}/\allowbreak{}p\allowbreak{}a\allowbreak{}i\allowbreak{}n\allowbreak{}t\allowbreak{}.\allowbreak{}t\allowbreak{}s} is intentionally simple because localization already happened upstream.

If paint ever needs to reason about negative roots, accumulated offsets, or tree rebasing, the ownership boundary has already been violated.

\section{The recursion model}\label{record-049-the-recursion-model}

\texttt{m\allowbreak{}e\allowbreak{}a\allowbreak{}s\allowbreak{}u\allowbreak{}r\allowbreak{}e\allowbreak{}L\allowbreak{}a\allowbreak{}y\allowbreak{}o\allowbreak{}u\allowbreak{}t\allowbreak{}B\allowbreak{}o\allowbreak{}u\allowbreak{}n\allowbreak{}d\allowbreak{}s\allowbreak{}(\allowbreak{}.\allowbreak{}.\allowbreak{}.\allowbreak{})} and \texttt{p\allowbreak{}a\allowbreak{}i\allowbreak{}n\allowbreak{}t\allowbreak{}L\allowbreak{}a\allowbreak{}y\allowbreak{}o\allowbreak{}u\allowbreak{}t\allowbreak{}N\allowbreak{}o\allowbreak{}d\allowbreak{}e\allowbreak{}W\allowbreak{}i\allowbreak{}t\allowbreak{}h\allowbreak{}O\allowbreak{}f\allowbreak{}f\allowbreak{}s\allowbreak{}e\allowbreak{}t\allowbreak{}(\allowbreak{}.\allowbreak{}.\allowbreak{}.\allowbreak{})} are both recursive, but they solve different problems.

\subsection{Bounds recursion}\label{record-049-bounds-recursion}

\texttt{m\allowbreak{}e\allowbreak{}a\allowbreak{}s\allowbreak{}u\allowbreak{}r\allowbreak{}e\allowbreak{}L\allowbreak{}a\allowbreak{}y\allowbreak{}o\allowbreak{}u\allowbreak{}t\allowbreak{}B\allowbreak{}o\allowbreak{}u\allowbreak{}n\allowbreak{}d\allowbreak{}s\allowbreak{}(\allowbreak{}.\allowbreak{}.\allowbreak{}.\allowbreak{})} walks the whole tree and accumulates:

\begin{itemize}
\tightlist
\item
  smallest \texttt{x}
\item
  smallest \texttt{y}
\item
  largest \texttt{x\allowbreak{} \allowbreak{}+\allowbreak{} \allowbreak{}w\allowbreak{}i\allowbreak{}d\allowbreak{}t\allowbreak{}h}
\item
  largest \texttt{y\allowbreak{} \allowbreak{}+\allowbreak{} \allowbreak{}h\allowbreak{}e\allowbreak{}i\allowbreak{}g\allowbreak{}h\allowbreak{}t}
\end{itemize}

This is the phase that determines the translation and final localized size.

\subsection{Paint recursion}\label{record-049-paint-recursion}

\texttt{p\allowbreak{}a\allowbreak{}i\allowbreak{}n\allowbreak{}t\allowbreak{}L\allowbreak{}a\allowbreak{}y\allowbreak{}o\allowbreak{}u\allowbreak{}t\allowbreak{}N\allowbreak{}o\allowbreak{}d\allowbreak{}e\allowbreak{}W\allowbreak{}i\allowbreak{}t\allowbreak{}h\allowbreak{}O\allowbreak{}f\allowbreak{}f\allowbreak{}s\allowbreak{}e\allowbreak{}t\allowbreak{}(\allowbreak{}.\allowbreak{}.\allowbreak{}.\allowbreak{})} or the paint middleware then walks the tree again and blits node surfaces into the already-sized target surface.

That recursion is intentionally dumb:

\begin{itemize}
\tightlist
\item
  no geometry solving
\item
  no bounds discovery
\item
  no re-localization
\end{itemize}

It is only applying the results of earlier layout ownership.

\section{When layout is recomputed today}\label{record-049-when-layout-is-recomputed-today}

In the base \texttt{r\allowbreak{}u\allowbreak{}n\allowbreak{}(\allowbreak{}.\allowbreak{}.\allowbreak{}.\allowbreak{})} runtime, layout localization is recomputed on every render.

That happens because the Layout stage does this every frame:

\begin{enumerate}
\def\labelenumi{\arabic{enumi}.}
\tightlist
\item
  call \texttt{a\allowbreak{}p\allowbreak{}p\allowbreak{}.\allowbreak{}v\allowbreak{}i\allowbreak{}e\allowbreak{}w\allowbreak{}(\allowbreak{}m\allowbreak{}o\allowbreak{}d\allowbreak{}e\allowbreak{}l\allowbreak{})}
\item
  call \texttt{w\allowbreak{}r\allowbreak{}a\allowbreak{}p\allowbreak{}V\allowbreak{}i\allowbreak{}e\allowbreak{}w\allowbreak{}O\allowbreak{}u\allowbreak{}t\allowbreak{}p\allowbreak{}u\allowbreak{}t\allowbreak{}A\allowbreak{}s\allowbreak{}L\allowbreak{}a\allowbreak{}y\allowbreak{}o\allowbreak{}u\allowbreak{}t\allowbreak{}R\allowbreak{}o\allowbreak{}o\allowbreak{}t\allowbreak{}(\allowbreak{}.\allowbreak{}.\allowbreak{}.\allowbreak{})}
\item
  store the result on \texttt{s\allowbreak{}t\allowbreak{}a\allowbreak{}t\allowbreak{}e\allowbreak{}.\allowbreak{}l\allowbreak{}a\allowbreak{}y\allowbreak{}o\allowbreak{}u\allowbreak{}t\allowbreak{}R\allowbreak{}o\allowbreak{}o\allowbreak{}t}
\end{enumerate}

So today there is no shared retained-layout cache inside the generic runtime render path.

Other paths follow the same basic rule:

\begin{itemize}
\tightlist
\item
  \texttt{n\allowbreak{}o\allowbreak{}r\allowbreak{}m\allowbreak{}a\allowbreak{}l\allowbreak{}i\allowbreak{}z\allowbreak{}e\allowbreak{}V\allowbreak{}i\allowbreak{}e\allowbreak{}w\allowbreak{}O\allowbreak{}u\allowbreak{}t\allowbreak{}p\allowbreak{}u\allowbreak{}t\allowbreak{}(\allowbreak{}.\allowbreak{}.\allowbreak{}.\allowbreak{})} recomputes localization when a scripted driver or test harness needs a \texttt{S\allowbreak{}u\allowbreak{}r\allowbreak{}f\allowbreak{}a\allowbreak{}c\allowbreak{}e}
\item
  \texttt{w\allowbreak{}r\allowbreak{}a\allowbreak{}p\allowbreak{}V\allowbreak{}i\allowbreak{}e\allowbreak{}w\allowbreak{}O\allowbreak{}u\allowbreak{}t\allowbreak{}p\allowbreak{}u\allowbreak{}t\allowbreak{}A\allowbreak{}s\allowbreak{}L\allowbreak{}a\allowbreak{}y\allowbreak{}o\allowbreak{}u\allowbreak{}t\allowbreak{}R\allowbreak{}o\allowbreak{}o\allowbreak{}t\allowbreak{}(\allowbreak{}.\allowbreak{}.\allowbreak{}.\allowbreak{})} recomputes localization when the pipeline needs a runtime-owned root tree
\end{itemize}

That is correct today, but it is not the final word on performance.

\section{Where retained layouts fit}\label{record-049-where-retained-layouts-fit}

Bijou already has retained-layout primitives in \texttt{r\allowbreak{}u\allowbreak{}n\allowbreak{}t\allowbreak{}i\allowbreak{}m\allowbreak{}e\allowbreak{}-\allowbreak{}e\allowbreak{}n\allowbreak{}g\allowbreak{}i\allowbreak{}n\allowbreak{}e\allowbreak{}.\allowbreak{}t\allowbreak{}s}, and the framed shell uses retained trees for input routing and hit testing.

The important current truth is:

\begin{itemize}
\tightlist
\item
  \textbf{retained layouts exist}
\item
  \textbf{the generic render pipeline does not yet retain and reuse \texttt{l\allowbreak{}a\allowbreak{}y\allowbreak{}o\allowbreak{}u\allowbreak{}t\allowbreak{}R\allowbreak{}o\allowbreak{}o\allowbreak{}t}}
\end{itemize}

So retained layout work should not be described as if the base runtime already paints from a persistent layout cache.

The likely future fit is:

\begin{enumerate}
\def\labelenumi{\arabic{enumi}.}
\tightlist
\item
  \texttt{v\allowbreak{}i\allowbreak{}e\allowbreak{}w\allowbreak{}(\allowbreak{}m\allowbreak{}o\allowbreak{}d\allowbreak{}e\allowbreak{}l\allowbreak{})} produces a logical tree in author space
\item
  localization establishes a runtime-owned local root
\item
  that localized root becomes the retained artifact
\item
  invalidation determines when it must be rebuilt
\item
  paint and hit testing both consume the same retained localized tree
\end{enumerate}

The key design pressure is consistency. A future retained-layout system should not make paint use one coordinate contract while hit testing uses another.

\section{Rule of thumb}\label{record-049-rule-of-thumb}

If you are working in this area:

\begin{itemize}
\tightlist
\item
  treat localization as the handoff from author coordinates to runtime coordinates
\item
  keep paint dumb
\item
  keep target-surface sizing separate from tree localization
\item
  describe retained layouts as a future reuse seam unless you are explicitly in the shell/runtime-engine retained-layout code
\end{itemize}

\section{Read this with}\label{record-049-read-this-with}

\begin{itemize}
\tightlist
\item
  \texttt{p\allowbreak{}a\allowbreak{}c\allowbreak{}k\allowbreak{}a\allowbreak{}g\allowbreak{}e\allowbreak{}s\allowbreak{}/\allowbreak{}b\allowbreak{}i\allowbreak{}j\allowbreak{}o\allowbreak{}u\allowbreak{}-\allowbreak{}t\allowbreak{}u\allowbreak{}i\allowbreak{}/\allowbreak{}s\allowbreak{}r\allowbreak{}c\allowbreak{}/\allowbreak{}v\allowbreak{}i\allowbreak{}e\allowbreak{}w\allowbreak{}-\allowbreak{}o\allowbreak{}u\allowbreak{}t\allowbreak{}p\allowbreak{}u\allowbreak{}t\allowbreak{}.\allowbreak{}t\allowbreak{}s}
\item
  \texttt{p\allowbreak{}a\allowbreak{}c\allowbreak{}k\allowbreak{}a\allowbreak{}g\allowbreak{}e\allowbreak{}s\allowbreak{}/\allowbreak{}b\allowbreak{}i\allowbreak{}j\allowbreak{}o\allowbreak{}u\allowbreak{}-\allowbreak{}t\allowbreak{}u\allowbreak{}i\allowbreak{}/\allowbreak{}s\allowbreak{}r\allowbreak{}c\allowbreak{}/\allowbreak{}l\allowbreak{}a\allowbreak{}y\allowbreak{}o\allowbreak{}u\allowbreak{}t\allowbreak{}-\allowbreak{}n\allowbreak{}o\allowbreak{}d\allowbreak{}e\allowbreak{}-\allowbreak{}s\allowbreak{}u\allowbreak{}r\allowbreak{}f\allowbreak{}a\allowbreak{}c\allowbreak{}e\allowbreak{}.\allowbreak{}t\allowbreak{}s}
\item
  \texttt{p\allowbreak{}a\allowbreak{}c\allowbreak{}k\allowbreak{}a\allowbreak{}g\allowbreak{}e\allowbreak{}s\allowbreak{}/\allowbreak{}b\allowbreak{}i\allowbreak{}j\allowbreak{}o\allowbreak{}u\allowbreak{}-\allowbreak{}t\allowbreak{}u\allowbreak{}i\allowbreak{}/\allowbreak{}s\allowbreak{}r\allowbreak{}c\allowbreak{}/\allowbreak{}r\allowbreak{}u\allowbreak{}n\allowbreak{}t\allowbreak{}i\allowbreak{}m\allowbreak{}e\allowbreak{}.\allowbreak{}t\allowbreak{}s}
\item
  \texttt{p\allowbreak{}a\allowbreak{}c\allowbreak{}k\allowbreak{}a\allowbreak{}g\allowbreak{}e\allowbreak{}s\allowbreak{}/\allowbreak{}b\allowbreak{}i\allowbreak{}j\allowbreak{}o\allowbreak{}u\allowbreak{}-\allowbreak{}t\allowbreak{}u\allowbreak{}i\allowbreak{}/\allowbreak{}s\allowbreak{}r\allowbreak{}c\allowbreak{}/\allowbreak{}p\allowbreak{}i\allowbreak{}p\allowbreak{}e\allowbreak{}l\allowbreak{}i\allowbreak{}n\allowbreak{}e\allowbreak{}/\allowbreak{}m\allowbreak{}i\allowbreak{}d\allowbreak{}d\allowbreak{}l\allowbreak{}e\allowbreak{}w\allowbreak{}a\allowbreak{}r\allowbreak{}e\allowbreak{}/\allowbreak{}p\allowbreak{}a\allowbreak{}i\allowbreak{}n\allowbreak{}t\allowbreak{}.\allowbreak{}t\allowbreak{}s}
\item
  \texttt{p\allowbreak{}a\allowbreak{}c\allowbreak{}k\allowbreak{}a\allowbreak{}g\allowbreak{}e\allowbreak{}s\allowbreak{}/\allowbreak{}b\allowbreak{}i\allowbreak{}j\allowbreak{}o\allowbreak{}u\allowbreak{}-\allowbreak{}t\allowbreak{}u\allowbreak{}i\allowbreak{}/\allowbreak{}s\allowbreak{}r\allowbreak{}c\allowbreak{}/\allowbreak{}r\allowbreak{}u\allowbreak{}n\allowbreak{}t\allowbreak{}i\allowbreak{}m\allowbreak{}e\allowbreak{}-\allowbreak{}e\allowbreak{}n\allowbreak{}g\allowbreak{}i\allowbreak{}n\allowbreak{}e\allowbreak{}.\allowbreak{}t\allowbreak{}s}
\item
  \texttt{p\allowbreak{}a\allowbreak{}c\allowbreak{}k\allowbreak{}a\allowbreak{}g\allowbreak{}e\allowbreak{}s\allowbreak{}/\allowbreak{}b\allowbreak{}i\allowbreak{}j\allowbreak{}o\allowbreak{}u\allowbreak{}-\allowbreak{}t\allowbreak{}u\allowbreak{}i\allowbreak{}/\allowbreak{}A\allowbreak{}R\allowbreak{}C\allowbreak{}H\allowbreak{}I\allowbreak{}T\allowbreak{}E\allowbreak{}C\allowbreak{}T\allowbreak{}U\allowbreak{}R\allowbreak{}E\allowbreak{}.\allowbreak{}m\allowbreak{}d}
\item
  \texttt{p\allowbreak{}a\allowbreak{}c\allowbreak{}k\allowbreak{}a\allowbreak{}g\allowbreak{}e\allowbreak{}s\allowbreak{}/\allowbreak{}b\allowbreak{}i\allowbreak{}j\allowbreak{}o\allowbreak{}u\allowbreak{}-\allowbreak{}t\allowbreak{}u\allowbreak{}i\allowbreak{}/\allowbreak{}A\allowbreak{}D\allowbreak{}V\allowbreak{}A\allowbreak{}N\allowbreak{}C\allowbreak{}E\allowbreak{}D\allowbreak{}\_\allowbreak{}G\allowbreak{}U\allowbreak{}I\allowbreak{}D\allowbreak{}E\allowbreak{}.\allowbreak{}m\allowbreak{}d}
\end{itemize}

